% Source: https://artofproblemsolving.com/community/c3h1405396_2016_mathcounts_nationals_team_round_8
% Stable identity: aops-topic-1405396-q08
The condition says that the digits of $\dfrac{1}{N}$ in places $10^{-3},\dots,10^{-11}$ are all $3$, that is, the fractional part of $\dfrac{100}{N}$ lies in $[0.333333333,\,0.333333334)$. Write $\dfrac{100}{N}=k+\dfrac{1}{3}+\varepsilon$ with an integer $k\ge0$ and $0\le\varepsilon<10^{-9}$. Then
\[
\varepsilon=\dfrac{100}{N}-k-\dfrac{1}{3}=\dfrac{300-(3k+1)N}{3N}.
\]
If $N>300$ then $\dfrac{100}{N}<\dfrac{1}{3}$ and the fractional part is too small, so $N\le300$. Then a positive numerator would give $\varepsilon\ge\dfrac{1}{3N}\ge\dfrac{1}{900}>10^{-9}$, so the numerator is $0$: $(3k+1)N=300$. Thus $\dfrac{300}{N}$ is a divisor of $300$ that is $\equiv1\pmod3$, namely $1,4,10,25,100$, giving
\[
N=300,\ 75,\ 30,\ 12,\ 3 .
\]
(For each, $\dfrac{1}{N}=0.00\overline{3},\ 0.01\overline{3},\ 0.0\overline{3},\ 0.08\overline{3},\ 0.\overline{3}$, and from the third decimal place on all digits are $3$.) The sum is $300+75+30+12+3=\boxed{420}$.
